\documentclass[10pt,twocolumn,letterpaper]{article}

\usepackage[final]{cvpr}
\usepackage{times}
\usepackage{graphicx}
\usepackage{amsmath}
\usepackage{amssymb}
\usepackage{booktabs}
\usepackage{tabularx}
\usepackage{array}

\def\paperID{0000}
\def\confName{CVPR}
\def\confYear{2027}

\title{Supplementary Material: Learning Reusable Geometry Programs for Fast Tight Bounding Box Fitting}

\author{Chanhyeok Park \and Minhyuk Sung}

\begin{document}
\maketitle

\section{Overview}

This supplementary document records implementation details for the learned
SMART~\cite{park2024smart} router and variable-length macro-skill experiments.
The key principle is unchanged from the main paper: learned components propose
actions, but exact SMART/Manifold~\cite{huang2020manifoldplus} reward validates
accepted updates.

\section{Token Schema}

\subsection{Shape Tokens}

Shape tokens summarize the current refine state:

\begin{itemize}
\item category id or category embedding,
\item current turn,
\item number of active boxes,
\item current exact score if available,
\item coverage and BVS proxy summaries,
\item centroid extent and PCA-like extent,
\item volume histogram or occupancy summary.
\end{itemize}

\subsection{Box Tokens}

Each active box contributes:

\begin{itemize}
\item center in normalized object coordinates,
\item side lengths,
\item volume and relative volume,
\item aspect ratio statistics,
\item rotation summary,
\item slack/coverage contribution proxies.
\end{itemize}

\subsection{Candidate Tokens}

Each candidate primitive edit contributes:

\begin{itemize}
\item action type,
\item affected box id,
\item affected face, axis, or center coordinate,
\item signed delta/scale,
\item proxy coverage after action,
\item proxy BVS after action,
\item proxy rank and z-score within the current candidate set,
\item invalidity and guard flags.
\end{itemize}

Offline training files additionally store exact SMART reward for every
candidate.  This exact reward is the supervision target, not an inference-time
replacement.

\section{Skill Representation}

The current mined skills are stored as parameterized programs.  A typical
record has the following logical form:

\begin{verbatim}
{
  "skill": "repeat_shrink_recenter",
  "precondition": {
    "category": ["airplane", "chair", "table"],
    "feature": "max_bbox_volume_ratio",
    "op": "<",
    "threshold": 1.5581414368932673
  },
  "program": [
    {"op": "shrink_face", "repeat": "upper_quantile"},
    {"op": "recenter_box", "repeat": 1},
    {"op": "shrink_face", "repeat": "upper_quantile"}
  ],
  "termination": {
    "max_steps": 16,
    "accept": "exact_non_worse"
  }
}
\end{verbatim}

The program is intentionally not a fixed raw action sequence.  It can be
instantiated with different boxes, faces, and repeat counts depending on the
state.  This is the difference between memorizing a trace and extracting a
reusable geometry skill.

\section{Observed Skill Families}

\paragraph{Shrink--recenter--shrink.}
This is the strongest current family.  It tends to appear when a box already
covers the relevant mass but has excess slack.  A recenter step between shrink
bursts prevents the box from drifting into a bad local configuration.

\paragraph{Coverage rescue then tighten.}
This family first expands or recenters toward uncovered mass, then shrinks
slack faces once coverage is recovered.  It is useful for local-minimum escape
because a temporary score-neutral or score-negative expansion can enable a
later score-positive tightening sequence.

\paragraph{Major-axis extend then trim.}
This pattern appears in elongated structures where a box must first align with
or extend along the object's dominant direction before minor faces can be
tightened safely.

\section{Conservative Repeat-Budget Gate}

The current strongest variable-length skill result uses a conservative gate for
longer repeat budgets.  The mined rule is:

\begin{verbatim}
airplane/chair -> upper-repeat
table -> upper-repeat only if num_actions >= 130
\end{verbatim}

The interpretation is that table states should use a longer repeat budget only
when the current decomposition exposes enough legal actions.  Low-action-count
one-box table states are unsafe because the longer repeat program can fail to
recover from under-decomposition.

\begin{table}[h]
\centering
\small
\begin{tabular}{lccc}
\toprule
Policy & States & Mean delta & Accepted \\
\midrule
Median repeat & 173 & $+1.5200$ & 173/173 \\
Blind upper repeat & 173 & $+1.5447$ & 172/173 \\
Conservative gate & 173 & $+1.5759$ & 173/173 \\
\bottomrule
\end{tabular}
\caption{Full hard-state validation for the repeat-budget gate.}
\end{table}

On table states alone, blind upper-repeat gives 10 wins, 28 ties, and 5 losses
against median-repeat, with negative mean difference.  The conservative gate
gives 7 wins, 36 ties, and 0 losses.  This supports the option view: the skill
is useful, but only under a geometric precondition.

\paragraph{Runtime versus mining mode.}
The runtime controller uses first-positive acceptance and averages 1.023 exact
skill attempts on the hard-state portfolio.  For analysis only, best-positive
acceptance exact-scores all top-3 attempts, improving mean delta from
$+1.5759$ to $+1.6020$ but increasing mean time from 0.177s to 0.380s.  The
extra evaluations reveal that hard states require learning option parameters:
repeat budget, target role, and whether the controller should use direct long
shrink or a recenter-then-shrink program.

\paragraph{Pattern-memory replay.}
On the same hard-state candidate logs, a macro-heavy memory score over
macro-id, canonical pattern, family, and skill recovers most of the logged
top-3 oracle.  One exact call improves the held-out replay mean from $+1.1207$
to $+1.4675$, while two exact calls reach $+1.5267$ versus the three-call
oracle at $+1.5269$.  This is a bounded replay estimate, but it shows that
reusable program identity is a stronger signal than semantic family alone.

\paragraph{Live hybrid selector.}
We also tested a live hybrid selector that first builds a top-3 pool with the
program gate, then reorders only that pool using macro-memory and exact-scores
the top two candidates.  On the hard-state portfolio this gives $+1.6017$ mean
delta with 2.0 exact attempts, compared with $+1.6020$ for top-3
best-positive at 3.0 exact attempts.  Thus the hybrid selector removes one
third of exact skill evaluations with negligible loss inside the current
candidate set.

\paragraph{Conditional second exact.}
The top-2 logs show that the second exact skill wins only $31/173$ states, with
$25$ of those wins on table states.  A single branch-opening precondition,
\texttt{category=table} and \texttt{num\_actions < 138}, opens the second exact
skill on $34/173$ states.  This gives $+1.6015$ mean delta, $1.197$ exact
attempts, and $0.195$s mean elapsed time, compared with $+1.6020$, $3.000$
attempts, and $0.380$s for top-3 best-positive.  This is a $60.1\%$ exact-call
reduction and a $1.95\times$ per-state speedup with $0.00044$ mean delta loss.
In five-fold train/test mining, the held-out gate gives $+1.5984$ mean delta
with $1.197$ exact attempts, compared with $+1.5996$ for the top-2 exact oracle
and $+1.5478$ for first-skill only.

\paragraph{Conditional exact-budget ladder.}
We also mined the exact budget as an option precondition.  Instead of always
checking top-3 skills, the controller chooses budget 1, 2, or 3 from cheap state
rules.  On the runnable 133-state portfolio, the mined rule is category-level:
airplane states use one exact skill check, chair states use two, and table
states use three.  Live execution gives $+1.7237$ mean delta with $1.669$ exact
attempts and $0.254$s mean elapsed time, compared with $+1.7247$, $3.000$
attempts, and $0.292$s for top-3 exact.  On the harder 173-state conservative
portfolio, a quality-preserving bbox-volume rule reaches $+1.60175$ versus
$+1.60196$ for top-3 exact, with $2.595$ exact attempts.  These rules are less
aggressive than the speed-first second-exact gate, but they preserve top-3
quality more tightly.

\section{Validation Checklist}

Before promoting any learned or macro path to a default setting, the following
checks should pass:

\begin{itemize}
\item accepted updates are exact-score non-worse than baseline under the target
budget,
\item rollback leaves the native C++ engine state unchanged after failed skill
attempts,
\item candidate pruning has zero losses on a frozen held-out split,
\item learned MCTS prior improves or ties score under equal or lower wall time,
\item macro preconditions are selected using validation data only, not test
data,
\item rare category-specific failures are reported separately.
\end{itemize}

\section{Benchmark Manifest Template}

For paper experiments, each benchmark run should be declared in a manifest
rather than assembled manually.  A minimal manifest should include:

\begin{verbatim}
experiment:
  name: guarded_mcts_prior_v1
  seed: 0
  engine: cpp_native
  exact_reward: manifold

data:
  train_states: path/to/train.jsonl
  val_states: path/to/val.jsonl
  test_states: path/to/test.jsonl
  categories: [airplane, chair, table]

baseline:
  refine_turns: 4
  mcts_iter: 50
  mcts_depth: 2

learned:
  router_checkpoint: path/to/router.smartmlp
  profile: guarded
  topk_multibox: 1
  topk_onebox: 15
  transposition_table: true

metrics:
  - exact_score_delta
  - exact_call_count
  - wall_time
  - quality_losses
  - rollback_count
\end{verbatim}

The test split should be evaluated once after model and precondition selection.
Any additional tuning after seeing test failures should create a new frozen
split.

\section{Current Research Caveat}

The macro-skill evidence is promising but not yet a release default.  The
largest current risk is overfitting preconditions to a small number of rare
failure examples.  The next data collection round should deliberately increase
hard table, chair-leg, thin-part, and disconnected-component cases.

{\small
\bibliographystyle{ieeenat_fullname}
\bibliography{refs}
}

\end{document}
